\section{Quorum-protected files}
\label{sec:quorum-files}

\code{access quorum} protects a file against a split tree the same way
\code{split} escrows a standalone secret --- \flag{--leaf} builds the tree
in the same lock command.

\begin{lstlisting}
keyquorum split --label "escrow demo" --threshold 2 \
  --leaf alice=alice.pub --leaf bob=bob.pub --generate-keys --register
keyquorum tree <key-id>
keyquorum reconstruct <key-id> --share-file alice.key --share-file bob.pub

keyquorum access quorum --state 0 --source ./secret.txt \
  --encrypted-path ./secret.txt.kqenc \
  --leaf alice=alice.pub --leaf bob=bob.pub --register
keyquorum access quorum --status --id <file-id>
keyquorum access quorum --state 1 --id <file-id> \
  --share-file alice.key --share-file bob.pub
\end{lstlisting}

\code{<key-id>} and \code{<file-id>} are printed by \code{split} / \code{access
quorum --state 0} (``Split key 1'' / ``Locked file 1''). \flag{--share-file}
takes the same key files \code{generate} wrote (\file{alice.pub} /
\file{alice.key}) or another standard public-key file (PEM, OpenSSH
\file{.pub}). The private key unwraps every leaf sealed to that hardware
key. \code{tree} still prints leaf node ids for inspection.

\subsection{Expiry}

\flag{--expires "yyyy-mm-dd hh:mm"} (UTC) gives the locked file a cutoff:
the first unlock attempt after that instant deletes the ciphertext and the
database row instead of unlocking it.

\begin{lstlisting}
keyquorum access quorum --state 0 --source ./secret.txt \
  --encrypted-path ./secret.txt.kqenc --leaf alice=alice.pub \
  --leaf bob=bob.pub --expires "2026-12-31 23:59"
\end{lstlisting}

\subsection{Verbose unlock}

\flag{--state 1 --verbose} prints which shares unwrapped, which subset met
the threshold, the physical devices counted toward
\code{minimum_physical_devices}, and any parent approval that was checked.
This is the most direct way to confirm \emph{why} an unlock succeeded or
failed on a given set of shares.
